\documentclass[12pt,a4paper]{article}
\usepackage{amsmath,amssymb,mathrsfs,tikz,times,pifont}
\usepackage{enumitem}
\usepackage{float}
\newcommand\circitem[1]{%
\tikz[baseline=(char.base)]{
\node[circle,draw=gray, fill=red!55,
minimum size=1.2em,inner sep=0] (char) {#1};}}
\newcommand\boxitem[1]{%
\tikz[baseline=(char.base)]{
\node[fill=cyan,
minimum size=1.2em,inner sep=0] (char) {#1};}}
\setlist[enumerate,1]{label=\protect\circitem{\arabic*}}
\setlist[enumerate,2]{label=\protect\boxitem{\alph*}}
%%%::::::by chnini ameur :::::::%%%
\everymath{\displaystyle}
\usepackage[left=1cm,right=1cm,top=1cm,bottom=1.7cm]{geometry}
\usepackage{array,multirow}
\usepackage[most]{tcolorbox}
\usepackage{varwidth}
\tcbuselibrary{skins,hooks}
\usetikzlibrary{patterns}
%%%::::::by chnini ameur :::::::%%%
\newtcolorbox{exa}[2][]{enhanced,breakable,before skip=2mm,after skip=5mm,
colback=yellow!20!white,colframe=black!20!blue,boxrule=0.5mm,
attach boxed title to top left ={xshift=0.6cm,yshift*=1mm-\tcboxedtitleheight},
fonttitle=\bfseries,
title={#2},#1,
% varwidth boxed title*=-3cm,
boxed title style={frame code={
\path[fill=tcbcolback!30!black]
([yshift=-1mm,xshift=-1mm]frame.north west)
arc[start angle=0,end angle=180,radius=1mm]
([yshift=-1mm,xshift=1mm]frame.north east)
arc[start angle=180,end angle=0,radius=1mm];
\path[left color=tcbcolback!60!black,right color = tcbcolback!60!black,
middle color = tcbcolback!80!black]
([xshift=-2mm]frame.north west) -- ([xshift=2mm]frame.north east)
[rounded corners=1mm]-- ([xshift=1mm,yshift=-1mm]frame.north east)
-- (frame.south east) -- (frame.south west)
-- ([xshift=-1mm,yshift=-1mm]frame.north west)
[sharp corners]-- cycle;
},interior engine=empty,
},interior style={top color=yellow!5}}
%%%%%%%%%%%%%%%%%%%%%%%

\usepackage{fancyhdr}
\usepackage{eso-pic}         % Pour ajouter des éléments en arrière-plan
% Commande pour ajouter du texte en arrière-plan
\AddToShipoutPicture{
    \AtTextCenter{%
        \makebox[0pt]{\rotatebox{80}{\textcolor[gray]{0.7}{\fontsize{5cm}{5cm}\selectfont PGB}}}
    }
}
\usepackage{lastpage}
\fancyhf{}
\pagestyle{fancy}
\renewcommand{\footrulewidth}{1pt}
\renewcommand{\headrulewidth}{0pt}
\renewcommand{\footruleskip}{10pt}
\fancyfoot[R]{
\color{blue}\ding{45}\ \textbf{2025}
}
\fancyfoot[L]{
\color{blue}\ding{45}\ \textbf{Prof:M. BA}
}
\cfoot{\bf
\thepage /
\pageref{LastPage}}
% Définition de l'encadré adaptatif avec fond jaune
\newtcolorbox{resultbox}{
    colback=red!30, % Fond rouge clair
    colframe=black, % Bordure noire fine
    sharp corners, % Coins nets
    boxrule=0.5pt, % Contour léger
    boxsep=2pt, % Espacement interne
    left=5pt, right=5pt, top=2pt, bottom=2pt, % Marges internes
}
\begin{document}
\renewcommand{\arraystretch}{1.5}
\renewcommand{\arrayrulewidth}{1.2pt}
\begin{tikzpicture}[overlay,remember picture]
    \node[draw=blue,line width=1.2pt,fill=purple,text=blue,inner sep=3mm,rounded corners,pattern=dots]at ([yshift=-2.5cm]current page.north) {\begingroup\setlength{\fboxsep}{0pt}\colorbox{white}{\begin{tabular}{|*1{>{\centering \arraybackslash}p{0.28\textwidth}} |*2{>{\centering \arraybackslash}p{0.2\textwidth}|} *1{>{\centering \arraybackslash}p{0.19\textwidth}|} }
                \hline
                \multicolumn{3}{|c|}{$\diamond$$\diamond$$\diamond$\ \textbf{Lycée de Dindéfélo}\ $\diamond$$\diamond$$\diamond$ } & \textbf{A.S. : 202/2025}                                                                     \\ \hline
                \textbf{Matière: Mathématiques}                                                                                    & \textbf{Niveau : 1}\textbf{$^{er}$S2} & \textbf{Date: 16/06/2025} & \textbf{Durée : 4 heures} \\ \hline
                \multicolumn{4}{|c|}{\parbox[c]{10cm}{\begin{center}
                                                                  \textbf{{\Large\sffamily Composition n$ ^{\circ} $ 2 Du 2$ ^\text{\bf nd} $ Semestre}}
                                                              \end{center}}}                                                                                                                               \\ \hline
            \end{tabular}}\endgroup};
\end{tikzpicture}
\vspace{3cm}

\section*{Correction : Exercice 1 – Barycentres et lieux géométriques}

\textbf{Donné :} $ABCD$ un parallélogramme et $G = \mathrm{bar}\{(A,3);(B,2);(C,3);(D,2)\}$

\begin{enumerate}
    \item \textbf{Construction des barycentres $E$ et $F$} :
    \[
    E = \mathrm{bar}\{(A,3);(B,2)\} \quad \implies \quad AE : EB = 2:3
    \]
    \[
    F = \mathrm{bar}\{(C,3);(D,2)\} \quad \implies \quad CF : FD = 2:3
    \]
    
    \item \textbf{Montrer que $G$ est le milieu de $[EF]$ :}
    \[
    G = \mathrm{bar}\{(E,5);(F,5)\} \implies G \text{ milieu de } [EF]
    \]
    
    \item \textbf{Points $I$ et $J$ :}
    \begin{enumerate}
        \item $I$ centre de gravité de $ABC$ :
        \[
        I = \mathrm{bar}\{(A,1);(B,1);(C,1)\}
        \]
        \item $J$ symétrique de $B$ par rapport à $D$ :
        \[
        J = 2D - B = \mathrm{bar}\{(B;-1);(D,2)\}
        \]
        \item $I, J, G$ alignés car :
        \[
        G = \mathrm{bar}\{(I,3);(J,2)\}
        \]
    \end{enumerate}
    
    \item \textbf{Lieux géométriques :}
    \begin{enumerate}
        \item Équation $\|3\overrightarrow{MA} + 2\overrightarrow{MB}\| = \|3\overrightarrow{MC} + 2\overrightarrow{MD}\|$ :
        \[
        3\overrightarrow{MA} + 2\overrightarrow{MB} = 5\overrightarrow{ME}, \quad 3\overrightarrow{MC} + 2\overrightarrow{MD} = 5\overrightarrow{MF}
        \]
        Donc $\|\overrightarrow{ME}\| = \|\overrightarrow{MF}\|$  
        $\implies$ \textbf{médiatrice de $[EF]$}.
        
        \item Équation $\|3\overrightarrow{MA} + 2\overrightarrow{MB} + 3\overrightarrow{MC} + 2\overrightarrow{MD}\| = 20$ :
        \[
        3\overrightarrow{MA} + 2\overrightarrow{MB} + 3\overrightarrow{MC} + 2\overrightarrow{MD} = 5 \overrightarrow{MG}
        \]
        \[
        5 \|\overrightarrow{MG}\| = 20 \implies \|\overrightarrow{MG}\| = 4
        \]
        $\implies$ \textbf{cercle de centre $G$ et rayon $4$}.
    \end{enumerate}
\end{enumerate}

\bigskip
\textbf{Schéma :}
\begin{comment}
\begin{center}
\begin{tikzpicture}[scale=1.2]
    % Points du parallélogramme
    \coordinate (A) at (0,0);
    \coordinate (B) at (4,0);
    \coordinate (D) at (1,3);
    \coordinate (C) at ($(B)+(D)-(A)$);
    
    % Barycentres partiels
    \coordinate (E) at ($(A)!2/5!(B)$);
    \coordinate (F) at ($(C)!2/5!(D)$);
    
    % Barycentre G
    \coordinate (G) at ($(E)!0.5!(F)$);
    
    % Centre de gravité I
    \coordinate (I) at ($(A)!1/3!(B)!1/3!(C)$);
    
    % Symétrique J
    \coordinate (J) at ($(D)!2!(B)$);
    
    % Tracé du parallélogramme
    \draw [thick] (A) -- (B) -- (C) -- (D) -- cycle;
    
    % Tracé des barycentres
    \fill[red] (E) circle (2pt) node[below left] {$E$};
    \fill[red] (F) circle (2pt) node[below right] {$F$};
    \fill[blue] (G) circle (2pt) node[above] {$G$};
    \fill[green!70!black] (I) circle (2pt) node[left] {$I$};
    \fill[orange] (J) circle (2pt) node[right] {$J$};
    
    % Tracé de la médiatrice de [EF]
    \draw[dashed] (E) -- (F);
    \draw[dotted] ($(E)!0.5!(F)+(-0.5,1)$) -- ($(E)!0.5!(F)+(0.5,-1)$);
    
    % Cercle de centre G et rayon 4
    \draw[thick,->,>=latex] (G) -- ++(2,2); % pour repérer le rayon
    \draw[blue, thick] (G) circle (4);
    
    % Tracé des sommets
    \foreach \p/\l in {A/A,B/B,C/C,D/D}
        \fill (\p) circle (2pt) node[below] {\l};
\end{tikzpicture}
\end{center}
\end{comment}
\end{document}